\chapter{Implementacja}
\label{cha:implementacja}

\section{Środowisko}

Projekt uruchamiano w środowisku Python~3.11. Główne biblioteki:
\begin{itemize}
    \item \textbf{numpy} --- operacje na wektorach i macierzach,
    \item \textbf{deepface}, \textbf{tf-keras} --- moduł twarzy,
    \item \textbf{librosa}, \textbf{soundfile} --- przetwarzanie audio,
    \item \textbf{pytest} --- testy jednostkowe.
\end{itemize}

Zalecane jest użycie środowiska wirtualnego (\texttt{venv}). Szczegółowa instrukcja uruchomienia znajduje się w pliku \texttt{README.md}.

\section{Struktura projektu}

\begin{verbatim}
MultiVerify/
+-- main.py                 # identyfikacja 1:N
+-- register.py             # rejestracja użytkowników
+-- src/
|   +-- modalities/
|   |   +-- base.py         # interfejs BiometricModality
|   |   +-- face/           # FaceModality (DeepFace)
|   |   +-- voice/          # VoiceModality (MFCC)
|   +-- fusion/
|       +-- fusion_engine.py
+-- scripts/run_experiments.py
+-- data/users/<user>/      # dane biometryczne (lokalnie)
+-- tests/
\end{verbatim}

\section{Interfejs modalności}

Każda modalność implementuje klasę bazową z metodami \texttt{extract\_features} oraz \texttt{verify}, co umożliwia jednolite traktowanie różnych źródeł biometrii.

\section{Silnik fuzji}

Listing~\ref{lst:fusion} przedstawia fragment implementacji strategii fuzji.

\lstinputlisting[style=pythonStyle, caption={Fragment silnika fuzji (\texttt{FusionEngine}).}, label=lst:fusion]{src/fusion_engine.py}

\section{Identyfikacja użytkownika}

Program \texttt{main.py} udostępnia interfejs CLI z parametrami \texttt{--face}, \texttt{--voice}, \texttt{--threshold} oraz \texttt{--strategy}. Fragment parsera argumentów przedstawiono w listingu~\ref{lst:main}.

\lstinputlisting[style=pythonStyle, caption={Parser argumentów w \texttt{main.py}.}, label=lst:main]{src/main_args.py}

\section{Rejestracja}

Skrypt \texttt{register.py} automatycznie przetwarza wszystkie podkatalogi w \texttt{data/users/}, generując pliki \texttt{.npy} z embeddingami. Dzięki temu identyfikacja nie wymaga ponownego przetwarzania surowych plików multimedialnych.

\section{Testy}

Zaimplementowano testy jednostkowe modułu fuzji oraz funkcji podobieństwa. Uruchomienie: \texttt{pytest}. W momencie opracowania dokumentacji wszystkie 11 testów przechodziło pomyślnie.
